% DELTA23, /root/characters.  Source overlap is stated explicitly below.

\subsection{An induced harmonic quotient and the exact finite ambiguity}

\paragraph{D23.C11 (proved; an audit of the proposed induced $S_4$ quotient).}
Induce the marked permutation representation
\[
 F\longrightarrow S_4,\qquad x\longmapsto(123),\quad
 y\longmapsto(1234),
\]
from the index-six pure subgroup to $B=C_2\amalg C_3$.
The induced representation has degree $24$.  Its $B$ image has
order $41472$ and its $F$ image has derived orders
\[
 6912,1728,64,1.
\]
In particular the abelianization of the first derived image is
a three-group of order $27$.

Join this representation to the signed projective congruence quotient
$Q_{12}$, represented faithfully on its $48$ primitive row vectors
modulo global sign.  The joint $B$ image has order $995328$ and
its pure $F$ image has order $165888$.  The centralizer of the
image of $s$ within that pure group has order $192$ and contains
an element $b$ with positive signed residue
\[
 q_{12}(b)=\begin{pmatrix}9&4\\8&9\end{pmatrix}.
\]
This residue is $I$ modulo four and $S$ modulo three, so its
harmonic parity is odd.  The same $b$ belongs to the first derived
pure image, of order $13824$, whose $s$-centralizer has order $96$.
Thus the finite conjugacy equation for a pure $s$-frame in this
quotient has an odd-parity ambiguity even with zero abelianized
frame.  This is a finite ambiguity assertion, not an assertion
that the ambiguous frame choices lift to actual gentle elements.

\emph{Proof.}
The companion script \path{agents/characters/code/induced_s4_harmonic.py}
uses the literal words
$x=(sr)^{-2}$ and $y=sxs$.  Its six coset representatives are
$1,s,r,r^2,sr,sr^2$.  Every Schreier cocycle is reduced in the
free product and rewritten exactly as a word in $x,y$; the full
rewriting table and the two induced permutations are recorded in
\path{agents/characters/records/induced_s4_harmonic.json}.
The checks $s^2=r^3=1$ and the restrictions
$x|_{\text{identity block}}=(123)$,
$y|_{\text{identity block}}=(1234)$ verify the source marking.
Exact permutation-group stabilizer chains give the displayed orders
and derived subgroups.  The joint image uses the same right-action
convention on primitive row vectors, so its generators are the
same literal $s,r$ in both factors.
For comparison with FP.37's left-action convention, inversion of
the two four-point permutations gives $(p^{-1},t^{-1})$.
This pair is simultaneously conjugate to $(p,t)$ by $(13)$;
thus it has exactly the same marked free-group kernel.  This
convention conversion does not alter the induced quotient used here.

The displayed centralizer witness has induced permutation
\[
 (0\ 1)(2\ 3)(4\ 5)(6\ 7)(12\ 15)(13\ 14)
             (16\ 18)(17\ 19),
\]
with the other points fixed, using zero-based labels.
Its full $72$-point permutation is in the certificate.  Exact
membership in the derived pure group and commutation with the
joint image of $s$ are checked.  The matrix calculation modulo
four and three gives the asserted parity.\hfill$\square$

The Ferrari extension of this quotient contributes no additional
pentagon equation by itself.  Under $a=c=u,b=v$ one has
$(X_1,X_2,X_3,X_4,X_5)=(x,y,x,1,1)$; the two surviving terms
are $f(y,x)f(x,y)=1$.  This is precisely the general forgetful
obstruction already proved in source FP.30, not a new result.

\subsection{A graph proof of the existing relative homology detector}

\paragraph{D23.C12 (proved; independent graph reproof and exact implementation).}
Let $p$ be prime, $q\ge2$, and
$\Gamma=C_p\amalg C_q=\langle a,b\rangle$ be the profinite
free product.  Let $\pi:\Gamma\twoheadrightarrow Q$ be a finite
quotient injective on both factors, and set $U=\ker\pi$.
Then
\[
 \operatorname{pr}_Q C_{\Gamma/([U,U]U^p)}(a)
       =\langle\pi(a)\rangle.
\]
For the literal quotient $Q_{12}$ of $B=C_2\amalg C_3$, the
kernel $U$ is free profinite of rank $97$.  The modulus-two
relative homology extension separates the harmonic frame residues,
and the modulus-three extension separates the cubic ones.  Their
combined quotient
\[
 \Gamma_6=B/([U,U]U^6),\qquad |\Gamma_6|=576\,6^{97},
\]
therefore determines both pure $Q_{12}$ frame residues from the
actual multiplier-one action, and hence determines $a_2,a_3$.

\emph{Proof.}
Use the Bass--Serre tree of the free product, with edges indexed
by $\Gamma$, and vertices indexed by the cosets of its two
factors.  The torsion-free kernel $U$ acts freely.  Its quotient
is the finite connected graph with edges $Q$ and vertices
$Q/\langle\pi(a)\rangle$ and $Q/\langle\pi(b)\rangle$.
Its cycle space over $\mathbb F_p$ identifies with
$U/[U,U]U^p$.  All statements may first be made for the discrete
free product and its finite-index free kernel, then completed;
the finite quotient and its finite homology are unchanged.

Choose as base vertex $o$ the identity coset of the first factor.
An element with $Q$-image $g$ has a path-chain coordinate $v$
with boundary $g o-o$.  Changing the lift by an element of $U$
adds an arbitrary cycle.  The generator $a$ fixes the base vertex
and has zero path-chain coordinate.  Consequently a lift of
$g\in C_Q(\pi(a))$ centralizes $a$ in the relative extension
only if some path chain $v$ with that boundary is fixed by $a$.
This follows directly from the path cocycle rule
$v_{gh}=v_g+g v_h$; the commutator condition is
$(a-1)v=0$.

The action of $a$ on the edge set $Q$ is free, with all orbits
of length $p$.  Every invariant edge chain is a linear combination
of their orbit sums.  At a fixed vertex, the $p$ contributions
of such an orbit have the same incidence sign, and sum to zero
over $\mathbb F_p$.  Hence the boundary of an invariant chain
has coefficient zero at every $a$-fixed vertex.

Both $o$ and $g o$ are $a$-fixed when $g$ centralizes $a$.
If they are distinct, the prescribed boundary $g o-o$ has a
nonzero coefficient at both, a contradiction.  Thus $g o=o$,
or $g\in\langle\pi(a)\rangle$.  Powers of $a$ themselves
centralize $a$ and project onto that factor, proving equality.

For $Q_{12}$, the quotient graph has $576$ edges and
$288+192$ vertices, hence rank $576-288-192+1=97$.
Apply the first assertion with $a=s,p=2$ and, after interchanging
the factors, with $a=r,p=3$.  The finite images of the two
factors meet the pure subgroup $q(F)$ trivially.  Therefore two
pure conjugators of $s$ with the same finite conjugate have the
same $Q_{12}$ residue; the analogous statement holds for $r$.
Actual frames supply existence.  Their positive signed residues
are selected by their identity reduction modulo four.  R3 then
recovers the shear from the frame abelianization modulo three,
and the corrected frame parities follow.\hfill$\square$

The independent computation
\path{agents/characters/code/check_relative_frame_homology.py}
constructs the actual graph, a spanning tree, and all $97$ cycle
coordinates.  It solves every lifting equation
$(a-1)m=-(a-1)v$ over the relevant field.  At two the original
centralizer has order $16$, its liftable image has order $2$,
and $\operatorname{rank}(s-1\mid H_1)=45$.
At three the corresponding numbers are $18,3,63$.
It also checks the fixed-vertex boundary obstruction directly
on every edge-orbit sum.  Full data are in
\path{agents/characters/records/relative_frame_homology.json}.

\paragraph{Exact overlap and remaining scope.}
The source already proves the same frame-separation mechanism:
FP.23, in \path{sources/DELTA_ON_GT/agent_five/five_point_results.tex},
uses $\mathbb F_2[Q/H_2]\oplus\mathbb F_3[Q/H_3]$ and fixed
cosets; RH.2, in \path{sources/DELTA_ON_GT/sections/A_homology.tex},
explicitly says that detector factors through the same $\Gamma_6$.
D23.C12 is a direct graph proof, its general $C_p*C_q$ formulation,
and an exact rank-$97$ implementation of the lifting obstruction.
It is not claimed as a new detector or a new pentagon constraint.
The rank-$129$ and rank-$4097$ Ferrari-kernel tests concern different
cover homologies and retain five-point information.  As FP.30 proves,
simply applying Ferrari and inducing any of the present four-point
detectors leaves the pentagon word identically one for every
two-cycle word.  A comparison with a representation that retains
the relevant five-point kernel remains necessary.

\subsection{A smaller quotient when only the two parities are needed}

\paragraph{D23.C13 (proved; strict compression of the supplied detector).}
There is an explicit quotient of the FP.23 source image with order
\[
 576\cdot2^{34}\cdot3^{52}
\]
such that the value of an actual multiplier-one gentle word $f$
in this quotient determines both corrected frame parities
$a_2,a_3$.  Its ambient permutation modules have dimensions
$72$ over $\mathbb F_2$ and $64$ over $\mathbb F_3$.
This is a quotient for reading actual normalized action values,
or equivalently the word value $f$; invariance of its kernel
under all gentle automorphisms is not asserted.  The stable
quotient $\Gamma_6$ of D23.C12 maps onto it, so an actual action
on $\Gamma_6$ supplies the required readout.

The original FP.23 source image has exact order
$576\cdot2^{97}\cdot3^{96}$.  Thus the present quotient is
strictly smaller; the proved readout here concerns the two
parities, with no full frame-residue claim.  It still gives no additional
pentagon equation after Ferrari extension alone.

\emph{Proof.}
Write $C_s=C_Q(q(s))$ and $C_r=C_Q(q(r))$, of orders $16$ and
$18$.  A signed lift of an element of $C_s$ commutes with $S$
at both primes: a possible simultaneous anticommutation is
excluded modulo four, since it would require a sum of two
squares to be $-1$ modulo four.  At three, such a lift is
$S^j$; at four it is $aI+bS$ with $a^2+b^2=1$.
Define
\[
 \eta_2(g)=j+b\pmod2.
\]
This is independent of the global sign of the lift: changing that
sign adds two to $j$ and leaves $b$ unchanged modulo two.
Both summands are homomorphisms, the second being reduction to
the two-element centralizer of $S$ over $\mathbb F_2$.
It is a nonzero character and $\eta_2(q(s))=1+1=0$.
On the positive pure centralizers with residue $I$ modulo four,
it is exactly the harmonic exponent parity.

For an element of $C_r$, choose a signed lift.  At both moduli
four and three its centralizer has the form
$\{\pm R_0^j: j\in\mathbb Z/3\}$.
The lift commutes with $R_0$ genuinely: anticommutation would
negate its trace $-1$, already impossible modulo three.
Let $e_4,e_3\in\mathbb F_2$ be the respective signs and put
\[
 \eta_3(g)=e_4+e_3.
\]
Changing the global lift changes both signs and preserves their
sum.  This is again a nonzero character, trivial on $q(r)$;
on a positive pure centralizer it is exactly the cubic sign.

Put $H_2^+=\ker\eta_2\subset C_s$ and
$H_3^+=\ker\eta_3\subset C_r$.  Their orders are $8,9$ and
they contain $q(s),q(r)$, respectively.  Form
\[
 V_2=\mathbb F_2[Q/H_2^+],\qquad
 V_3=\mathbb F_3[Q/H_3^+],
\]
of dimensions $72$ and $64$, and let $e_2,e_3$ denote the
identity-coset basis vectors.  As in FP.23 the elements
\[
 S^+=(e_2,0,q(s)),\qquad R^+=(0,e_3,q(r))
\]
in $(V_2\oplus V_3)\rtimes Q$ have orders two and three.
They define a finite homomorphism $\theta^+:B\to E^+$.
Since $\langle q(s)\rangle\subset H_2^+$ and
$\langle q(r)\rangle\subset H_3^+$, the two coset projections
give a surjective module map from FP.23's ambient modules to
these modules, compatible with both generators.  In particular
$E^+$ is a quotient of the actual FP.23 image.

If $(m,g)$ centralizes $S^+$, its projection centralizes $q(s)$
and satisfies $(1-q(s))m_2=(1-g)e_2$.  Passing to coinvariants
shows that the two $q(s)$-fixed cosets $H_2^+$ and $gH_2^+$
must coincide; hence $g\in H_2^+$.  Thus every centralizer
ambiguity has $\eta_2(g)=0$.  The same argument with $R^+$
gives $g\in H_3^+$ and $\eta_3(g)=0$.

Given the value $\theta^+(f)$ for an actual gentle element,
the normalized formulas give
\[
 \theta^+\Phi(s)=S^+\theta^+(f),\qquad
 \theta^+\Phi(r)=\theta^+(f)^{-1}R^+.
\]
Solve the finite conjugacy equations for frames in $\theta^+(F)$
whose positive signed $Q$-residues are $I$ modulo four.
Actual frames guarantee solutions.  The ratio of two harmonic
solutions centralizes $S^+$ and has positive residue $I$
modulo four, so its $\eta_2$ is precisely its possible change
of harmonic parity; that change is zero.  The corresponding
argument with $\eta_3$ proves uniqueness of the cubic parity.
The shear is also known: the positive signed value of $f$
has $I$ residue modulo four, and its upper-right entry modulo
three is $\rho$ by R3.  Hence $V=U^{8\rho}$ modulo twelve
is determined before the corrected parities are read.
This proves the asserted readout without assuming every finite
admissible conjugacy solution lifts to an actual gentle element.

For the exact size, the implementation uses $576$ Schreier
representatives for $U=\ker q$.  If $v_g$ is the affine
translation of a chosen representative and $t_a$ that of a
generator $a\in\{s,r\}$, the Schreier kernel generator has
translation
\[
 v_g+g t_a-v_{ga}.
\]
These $1152$ vectors span exactly the image of $U$ in each
module.  Exact row reduction gives ranks $34$ and $52$.
The joint image of $U$ is the full product of these elementary
abelian groups, because their orders are coprime and the joint
image surjects onto both.  Surjectivity onto $Q$ gives the
stated total order.  The same calculation in FP.23's original
modules gives ranks $97$ and $96$, proving strictness.
Every character equation, every coset obstruction, the Schreier
spans, and the image ranks are checked by
\path{agents/characters/code/check_compressed_parity_detector.py},
with exact records in
\path{agents/characters/records/compressed_parity_detector.json}.
Finally FP.30 applies to every quotient factoring through the
forgetful map, so compression does not change its automatic
pentagon conclusion.\hfill$\square$
